\section{Proof of the uniform Iitaka theorem}
\label{sec:main-proof}

\begin{proof}[Proof of Theorem~\ref{thm:main}]
First we deduce nonvanishing from pseudo-effectivity. The companion
good-model theorem, with the \(\Q\)-factorial dlt choice in its
construction \cite[Theorem~11.1 and Proposition~2.5]{OpenAIAbundance},
gives a log minimal model \((Y,\Delta_Y)\) of \((X,\Delta)\).
Its boundary is the strict transform of \(\Delta\) together with any
extracted prime divisors of coefficient one. Hence it is rational,
and \(D_Y=K_Y+\Delta_Y\) is nef and \(\Q\)-Cartier. The companion's
rational abundance theorem \cite[Theorem~1.1]{OpenAIAbundance} makes
\(D_Y\) semiample. On a common smooth resolution
\[
 X\xleftarrow{p}W\xrightarrow{q}Y
\]
with compatible canonical divisors, its comparison lemma
\cite[Lemma~2.2]{OpenAIAbundance} gives
\[
 p^*D=q^*D_Y+E,\qquad E\geq0,\qquad E\text{ is }q\text{-exceptional}.
\]
Choose a positive integer \(r\) such that \(rD_Y\) is Cartier and
globally generated, and a nonzero section
\(s\in H^0(Y,\OO_Y(rD_Y))\), viewed in the common rational function field.
Pulling the effective Cartier divisor
\(\Div_Y(s)+rD_Y\) back by \(q\), adding \(rE\), and pushing forward
by \(p\) gives
\[
 \Div_X(s)+rD\geq0.
\]
The pole test \eqref{eq:floor-pole-test} therefore gives
\(0\ne s\in V_r(D)\). Choose a positive integer \(t\) such that
\(trD\) is integral Cartier. Then \(s^t\) is a nonzero section of
\(\OO_X(trD)\), so \(\kappa(X,D)\geq0\).
These auxiliary degrees need not be uniform: they only establish the
hypothesis for the uniform argument, whose all-degree comparison
will retain the full Iitaka field.

Enlarge the fixed coefficient set to $\Phi'=\Phi\cup\{1\}$.
Proposition~\ref{prop:semiample-model} gives a semiample dlt model
$(X',\Delta')$ with coefficients in $\Phi'$ and identical rounded
section spaces in every integer degree. Write its contraction as
$f:X'\to Z$ and its adjoint as
\[
 K_{X'}+\Delta'\sim_{\Q}f^*A,
\]
where $A$ is ample rational Cartier and
$\dim Z=\kappa(X,K_X+\Delta)$.

If $\dim Z>0$, apply Proposition~\ref{prop:effective-base} to
$f:(X',\Delta')\to Z$, with total-space dimension four and fixed
coefficient set $\Phi'$. Since $A$ is ample, it is big. The proposition
gives one degree depending only on $\Phi'$ whose every positive
multiple is nonempty and generates exactly $\C(Z)=K(D)$ inside the
common rational function field. This same degree covers all positive
base dimensions. When $\dim Z=4$, the contraction is birational and
its field is $\C(X')$.

Suppose instead that $Z$ is a point. The adjoint of $(X',\Delta')$
is then \(\Q\)-linearly trivial. There are three exhaustive cases.
If the pair is not klt, the four-dimensional non-klt index theorem
\cite[Corollary~1.7]{JiangLiu2021} gives a common linearly trivial
multiple depending only on $\Phi'$. If it is klt and
$\Delta'\ne0$, use \cite[Theorem~1.14]{Xu2019}. Its hyperstandard
coefficient convention includes any prescribed finite rational set:
one may take the defining set $\{1-b:b\in\Phi'\}$ and denominator
one. If it is klt and $\Delta'=0$, use
Theorem~\ref{thm:fourfold-index-zero}. A bound for the positive
torsion order gives a common trivial multiple by taking the least
common multiple of the finitely many possible orders.

Take a common multiple of these zero-dimensional degrees and of the
degree obtained for positive-dimensional bases. The choice depends
only on $\Phi$ and dimension four. In dimension zero it yields a
nonempty system with image a point. In positive dimension,
the every-multiple clause of Proposition~\ref{prop:effective-base}
retains the whole Iitaka field.
Finally \eqref{eq:all-rounded-sections} transfers the complete
systems back to the original pair. This proves every assertion,
including the rounded reflexive convention and recovery of the
full Iitaka base field. If a canonical representative is replaced by
$K_X+\Div(h)$, then the floor divisor changes by $m\Div(h)$ in
every integer degree. Multiplication by $h^{-m}$ identifies the two
section spaces and leaves all ratios unchanged. The same integer
therefore works for all canonical representatives.
\end{proof}
